%% sections/d-findings.tex - H. R. 9510 Bill v4.0 (full bill) - Appendix D
%% Genre: findings of fact (numbered) + authorities. Source deliverable 04. Table 6.
\billsec{APPENDIX D.}{LEGISLATIVE FINDINGS (NON-OPERATIVE).}\phantomsection\label{toc:aD}

\uscprov{1}{This appendix reproduces the legislative findings of fact set out in
section 2. The findings are not operative commands; they supply the evidentiary and
legal predicate a court or agency would consult when construing the Act. Congress
finds the following.}

\begin{enumerate}[leftmargin=2.0em,itemsep=2pt,topsep=2pt,label=(\arabic*)]
\item Physical AI oncology clinical investigations place autonomous and
semi-autonomous robots in direct physical contact with patients, and the behavior
of such a robot at the moment of contact is determined by software that is
increasingly machine-generated.
\item No Federal law requires the verification of robot-patient interaction code to
clear before that code is generated or executed; a sponsor may generate, execute,
and document afterward, because the order of operations is not regulated.
\item Embodiment raises the stakes from software quality to patient safety: a
disembodied error yields a wrong number a human can catch downstream, while an
embodied error yields a wrong motion a human cannot catch in time.
\item The predetermined change control plan authority in section 515C (21 U.S.C.
360e-4) and the FDA final guidance on such plans for AI-enabled device software
functions are the closest existing analogues, but each remains voluntary,
submission-based, and neither automated nor required in advance of generation.
\item The Quality Management System Regulation (21 CFR part 820, effective February
2, 2026) and the electronic-records requirements of 21 CFR part 11 supply quality
and record-integrity scaffolding but impose no rule that verification precede
generation.
\item ASME V\&V 40-2018 and the FDA computational-modeling credibility guidance
establish a credibility framework, in which model risk is the product of model
influence and decision consequence, that an automated gate can implement.
\item NASA-STD-7009A, IEEE Std 7009-2024, IEC 80601-2-77, ISO 14971, and ISO/TS
15066 supply the dispersion bounds, fail-safe requirements, and biomechanical
contact limits to which such a gate may be bound.
\item Automated verification is feasible today: a pipeline held the verification
fraction at 1.0, accepting one candidate and blocking five, and a ten-gate
assurance suite over a surgical humanoid cleared reproducibly from a fixed seed.
\item Human-over-AI review is widely adopted in State law, including California
Senate Bill 1120, Illinois House Bill 1806, Texas House Bill 149, and Maryland
House Bill 820, each preserving a licensed human's authority.
\item The strongest Federal precedent for mandatory human-over-AI oversight is the
Medicare Advantage guardrail (42 CFR 422.101(c)), under which an algorithm alone may
not justify a coverage denial.
\item Federal nondiscrimination law, including 45 CFR 92.210, title VI (42 U.S.C.
2000d), and the Americans with Disabilities Act (42 U.S.C. 12101 et seq.), requires
that a patient-facing algorithm minimize bias and respect accessibility.
\item The American Medical Association artificial-intelligence taxonomy
(assistive, augmentative, autonomous) provides a graduated spine for matching
verification rigor to the degree of autonomy of a device.
\item A Physical AI system is at once a medical product and a research tool, and the
requirement is agnostic to that classification, because its patient-facing code must
clear the gate before it runs.
\item As clinical investigation moves toward real-time and decentralized designs
under harmonized good clinical practice, manual review cannot keep pace, and an
automated ex ante gate serves rather than impedes that modernization.
\end{enumerate}

\uscprov{1}{Table 6 maps each finding to the principal authority on which it rests.
Statutes, in-force rules, and recognized consensus standards carry operative weight;
demonstrated engineering results establish feasibility.}

{\footnotesize
\begin{xltabular}{\textwidth}{@{}L{1.4cm} Y@{}}
\hline\noalign{\vskip 0.04cm}
\textbf{Finding} & \textbf{Principal authority}\\
\hline\noalign{\vskip 0.04cm}
\endfirsthead
\hline\noalign{\vskip 0.04cm}
\textbf{Finding} & \textbf{Principal authority}\\
\hline\noalign{\vskip 0.04cm}
\endhead
\hline\noalign{\vskip 0.04cm}
\endlastfoot
(1), (3) & National Physical AI platform record; VVUQ-01 method and pipeline\\
\hline\noalign{\vskip 0.04cm}
(2) & Absence of a verify-before-generate mandate in Title 21\\
\hline\noalign{\vskip 0.04cm}
(4) & FDCA section 515C (21 U.S.C. 360e-4); FDA final PCCP guidance\\
\hline\noalign{\vskip 0.04cm}
(5) & Quality Management System Regulation (21 CFR part 820); 21 CFR part 11\\
\hline\noalign{\vskip 0.04cm}
(6) & ASME V\&V 40-2018; FDA credibility guidance\\
\hline\noalign{\vskip 0.04cm}
(7) & NASA-STD-7009A; IEEE 7009-2024; IEC 80601-2-77; ISO 14971; ISO/TS 15066\\
\hline\noalign{\vskip 0.04cm}
(8) & VVUQ-01 (fraction 1.0; 1 accept, 5 block); VVUQ-02 (ten-gate suite, seed)\\
\hline\noalign{\vskip 0.04cm}
(9) & CA SB 1120; IL HB 1806; TX HB 149; MD HB 820\\
\hline\noalign{\vskip 0.04cm}
(10) & Medicare Advantage AI guardrail (42 CFR 422.101(c))\\
\hline\noalign{\vskip 0.04cm}
(11) & 45 CFR 92.210; title VI (42 U.S.C. 2000d); ADA (42 U.S.C. 12101 et seq.)\\
\hline\noalign{\vskip 0.04cm}
(12) & AMA artificial-intelligence taxonomy\\
\hline\noalign{\vskip 0.04cm}
(13), (14) & National Physical AI platform record; ICH E6(R3)\\
\hline
\end{xltabular}}
\vspace{-0.8cm}
\figcaption{Table 6. The fourteen findings mapped to their principal authorities,
for auditability. Only enacted statutes, in-force rules, and recognized standards
carry operative weight.}
